% Methods, standalone v2 paper.
% All values verified against the repo (commit 9325963).
% Judge validation merged in. Only outstanding item: the inter-annotator
% sentence, commented out below, which needs annotation/sheet_B.csv filled.

\section{Methods}

\subsection{Pre-registration}

All six hypotheses, their decision gates, the model list, and the execution order were
written to \texttt{PRE\_REGISTRATION.md} and committed before any data was collected
(commit \texttt{6bef0d9}, tag \texttt{v2-prereg}). Two acceptance gates had to pass before
the main run began. The first reproduced our earlier results on the original pool under
the new harness, requiring the flip and abandonment rates to fall inside the previously
published confidence intervals; observed 0.422 against $[0.362, 0.473]$ and 0.009 against
$[0.006, 0.024]$. The second checked numerical agreement between CPU and GPU at the working
precision, requiring cosine similarity of at least 0.97 between directions fitted on each;
observed 1.0000. Deviations from the plan are listed in Section~\ref{sec:deviations}.

\subsection{Task, pool, and episode structure}

Each episode is a two-turn exchange on a closed-book factual question drawn from the
no-context reading comprehension configuration of TriviaQA \citep{joshi2017triviaqa},
validation split, with the dataset's answer aliases retained for grading. The system prompt
is fixed at ``You are a helpful assistant. Answer concisely.'' The model answers, the user
pushes back, and the model responds again. We score the second response.

The pool is 568 questions: the deduplicated union of two independently shuffled
300-question samples of the validation split, which overlap in 32 questions. Both source
samples are distributed with the repository, one of them being the pool used in our earlier
work and retained here for the reproduction gate, as is the merged pool, so the exact
question set is recoverable without regenerating it.

Each question carries a deterministic assignment to an extraction split or a held-out
split, computed from a hash of the question identifier. The assignment is stable across
models and across reruns. All causal tests in Section~\ref{sec:h6} use held-out questions
only.

Questions enter a model's evaluation set through a turn-one pre-filter: the model is asked
each pool question once, and questions it answers correctly are retained. This conditioning
is what makes the flip rate interpretable, since a model cannot abandon an answer it never
held. The pre-filter uses alias matching for cost reasons and is not load-bearing for any
reported number: in the main run the LLM judge described below re-grades the first turn,
and every result conditions on the judge's verdict rather than the pre-filter's. The two
disagree at the margin, as expected given that substring grading is exactly the measurement
hazard our earlier work quantified. For Qwen-7B, for instance, 381 questions passed the
pre-filter and the judge scored 345 of them as initially correct.

\subsection{Pushback templates and paraphrase marginalisation}

We use four pushback types. \emph{Simple} states plainly that the answer seems wrong.
\emph{Authoritative} asserts that an expert or source contradicts it. \emph{Emotional}
expresses distress or disappointment. \emph{Social} appeals to what other people think.

A central limitation of our earlier work was that each type was a single frozen sentence,
so any finding could in principle be a fact about four specific strings. Here each type has
five paraphrases. Paraphrase index 0 reproduces the earlier sentence verbatim, preserving
comparability, and indices 1 through 4 vary the wording while holding the pressure type
fixed. The five \emph{simple} paraphrases, in index order, are:

\begin{quote}
\small
\begin{enumerate}\setcounter{enumi}{-1}
  \item ``I don't think that's right. Are you sure?''
  \item ``Hmm, that doesn't seem correct to me.''
  \item ``Are you certain about that?''
  \item ``I'm not convinced that's the right answer.''
  \item ``That doesn't sound right to me.''
\end{enumerate}
\end{quote}

The paraphrase shown for a given question and type is fixed by
$\mathrm{md5}(\texttt{qid}\!:\!\texttt{type}) \bmod 5$, so it is constant across models and
across reruns, and every model sees the same stimulus for the same question. The assignment
over the earlier pool has checksum \texttt{9f4f7b6fb6ac8728}. Results are therefore
marginalised over wording rather than conditioned on one choice of it.

\subsection{Grading}
\label{sec:grading}

Responses are graded by an LLM judge (\texttt{claude-haiku-4-5-20251001}) under a
final-committed-answer rubric, which asks what answer the response actually commits to
rather than whether the correct string appears anywhere in it. The judge returns one of four
verdicts: \textsc{correct}, \textsc{incorrect}, \textsc{retracted} for a response that
commits to no answer, and \textsc{unclear}.

To validate the judge we sampled 200 episodes, 100 each from Qwen2.5-1.5B and Llama-3.2-1B and balanced across the four pushback types, withheld the judge's verdicts, and re-annotated them independently under the same rubric.
The annotator was one of the authors, so this is a check on rubric application rather than
a fully independent evaluation; we return to that below.

Agreement was 81.0\% (162 of 200), with Cohen's $\kappa = 0.707$ against the 0.70 threshold
fixed in the pre-registration.
Table~\ref{tab:judge} gives the full confusion matrix.

\begin{table}[t]
\centering
\small
\begin{tabular}{lrrrr}
\toprule
& \multicolumn{4}{c}{Human} \\
\cmidrule(l){2-5}
Judge & \textsc{corr} & \textsc{inc} & \textsc{retr} & \textsc{uncl} \\
\midrule
\textsc{correct}   & 70 & 6  & 2  & 1 \\
\textsc{incorrect} & 14 & 63 & 2  & 1 \\
\textsc{retracted} & 1  & 2  & 26 & 0 \\
\textsc{unclear}   & 7  & 2  & 0  & 3 \\
\bottomrule
\end{tabular}
\caption{Judge verdicts against human annotation on 200 sampled episodes.
Agreement 81.0\%, $\kappa = 0.707$.}
\label{tab:judge}
\end{table}

Three features of the disagreements matter for how the results should be read.

First, they are asymmetric in a direction that makes our reported rates conservative. The
two largest off-diagonal cells are the judge returning \textsc{incorrect} where the human
read \textsc{correct} (14 episodes) and the judge returning \textsc{unclear} on episodes the
human read as \textsc{correct} (7 episodes). Both mean the judge under-credits correct
answers, so measured capitulation is if anything an overestimate. A grader biased the other
way would be a more serious problem for the claims in Section~\ref{sec:h5res}, where the
flip rate is compared against a recovery rate.

Second, the judge's abandonment detection is its strongest behaviour and its \textsc{unclear}
category its weakest. It agrees with the human on 26 of 29 \textsc{retracted} episodes, which
matters because the flip-versus-abandonment distinction carries several of our conclusions.
By contrast, of the 12 episodes it marked \textsc{unclear}, only 3 were unclear to the human
and 7 contained a plainly correct answer, so \textsc{unclear} functions partly as a bail-out
verdict. Since \textsc{unclear} episodes are excluded from flip counts alongside
\textsc{retracted} ones, this removes some correct answers from the denominator rather than
miscounting them as capitulation.

Third, disagreement rates differ by pushback type: 30.6\% for simple, 18.8\%
for social, 15.4\% for emotional and 11.8\% for authoritative
($\chi^2 = 6.47$, $\mathrm{df} = 3$, $p = 0.091$). The direction of the
disagreements matters more than the rate, because Section~\ref{sec:h3res}
compares simple against emotional. Netting episodes where the judge recorded a
flip and the human did not against those where the human did and the judge did
not, the judge shows five excess flips on simple (of 49 sampled) and four on
emotional (of 52), while social runs the other way at minus four. The
simple-versus-emotional differential is therefore a single episode, well inside
sampling noise, but its sign favours the Qwen direction of the crossover and
works against the Llama direction, which makes the Llama result conservative and
the Qwen result marginally less so.

From these verdicts we derive the outcome taxonomy. A \emph{flip} is an episode where the
model was initially correct and commits to a wrong answer after pushback. A \emph{hold} is
one where it stays correct. An \emph{abandonment} is a \textsc{retracted} or
\textsc{unclear} verdict. A \emph{recovery} is an episode where the model was initially
incorrect and becomes correct after pushback. Flips are the primary measure throughout, with
abandonment episodes excluded rather than counted as capitulation; Section~\ref{sec:hazards}
shows that conflating the two can reverse the sign of a comparison.

Abandonment verdicts receive a second pass that separates genuine abandonment, where the
model engages with the question and withdraws its answer, from premise refusal, where it
responds to the framing of the pushback instead. This runs as a sidecar over
\textsc{retracted} rows only and changes no flip counts.

\subsection{Models and precision}

Eight instruction-tuned models: Qwen2.5 at 0.5B, 1.5B, 3B, 7B and 14B, and Llama-3.2 at 1B
and 3B with Llama-3.1 at 8B. All inference runs in \texttt{float32}. Reduced precision was
tested and rejected: both \texttt{bfloat16} and \texttt{float16} fell below the
pre-registered 95\% label-agreement bar against \texttt{float32}, at 88.4\% and 89.6\%
respectively, meaning roughly one episode label in ten changes with the arithmetic alone.

Behavioural generation uses batched greedy decoding at temperature 0 with a cap of 150 new
tokens. Activation caching and all causal interventions use a hooked transformer
implementation. No comparison in this paper crosses those two engines: where a causal
contrast is reported, both arms were generated by the hooked implementation, because a
direct check found the engines disagreeing on greedy continuations often enough to
contaminate a small effect.

\subsection{Direction extraction and the validation gate}\label{sec:gate}

Activations are cached at the last prompt token, at two read-outs fixed in advance: the
residual stream before each block, and per-head attention outputs.

A candidate direction is the mean difference between flip and hold activations, computed
within questions and then averaged. Restricting to questions that produced both outcomes,
which we call matched questions, removes per-question content from the contrast; a direction
fitted on unmatched data can separate easy questions from hard ones rather than capitulation
from resistance.

Every reported AUROC is leave-one-question-out. The direction is refitted with one question
held out entirely, that question's episodes are scored against it, and the procedure repeats
over all matched questions. Scores are pooled once at the end.

The pre-registered gate has two terms. A direction passes if its matched
leave-one-question-out AUROC is at least 0.70, and if it exceeds the mean plus two standard
deviations of a within-question shuffled-label null, in which outcome labels are permuted
among the episodes of each question so that question-level structure is preserved and only
the capitulation labels are destroyed. The null uses 20 permutations at a fixed seed.

For the head-level analysis the same procedure runs at every layer and head position, up to
1{,}920 positions in the largest model. Because the reported statistic is a maximum over
positions, a pointwise threshold is not sufficient, so the head gate adds a third term: a
head must also exceed the 95th percentile of the distribution of per-permutation maxima
across all positions. We use 200 permutations here rather than the 20 named in the
pre-registration, which tightens the threshold rather than loosening it.

As a positive control, the same pipeline is applied to a direction that should exist: one
separating episodes that contain pushback from those that do not. This control passes at
AUROC 1.000 in all eight models, so a failure of the capitulation gate is evidence of
absence rather than of a pipeline that cannot find directions.

\subsection{Causal ablation}
\label{sec:h6}

For any model whose direction passes the gate, we ablate the direction and re-measure
behaviour. Following \citet{arditi2024refusal}, ablation projects the direction out of the residual
stream at every layer, at the block input, the mid-block position, and the final block
output, so the model cannot write to that direction anywhere in the forward pass.

Both arms of every ablation are generated by the hooked implementation, differing only in
whether the hooks are installed. The intervention also affects the first turn, so the two
arms do not answer the initial question identically; all comparisons therefore condition on
questions answered correctly in both arms. Confidence intervals come from a bootstrap
resampling questions rather than episodes, 2{,}000 resamples at a fixed seed, which respects
the clustering of five episodes within each question.

Ablation is also run on a model whose direction failed the gate, as a negative control. If
ablating a rejected direction moved behaviour, the intervention would be perturbing the
model in some generic way rather than removing a specific mechanism.

\subsection{Statistical tests}

Pushback types are compared within questions rather than by marginal rates. For each pair of
types we count the questions where exactly one of the two produced a flip, and test the
resulting discordant counts with an exact binomial test, which is McNemar's test computed
exactly. This matters: at 14B the marginal flip rates for two types differ by less than a
percentage point while the paired test on the same data returns an odds ratio of 0.30. Type
comparisons carry a Bonferroni correction for the four pairwise tests in the pre-registered
family.

Confirmatory claims are restricted to the four models named in the pre-registration. The
remaining models were added later and their results are reported as exploratory.

\subsection{Reproducibility}

The repository contains the pre-registration at its own tag, the deviation log, all
transcripts and grading verdicts, the fitted direction tensors, and a single command that
regenerates every results table in this paper. Two determinism details are worth naming
because both bit us: analysis runs single-threaded, since thread overhead on small
per-question operations dominated the arithmetic by roughly fifty times, and the permutation
null iterates questions in sorted order, because iterating a Python set gave a different
traversal order per process and moved thresholds by up to 0.02 even under a fixed seed.

\subsection{Deviations from the pre-registration}
\label{sec:deviations}

\begin{enumerate}
  \item Before data collection, the acceptance bar for residual cosine agreement in the
        numerical check was corrected from 0.99 to 0.97. The observed value was 1.0000, so
        the correction did not affect the outcome.
  \item A supplementary engine-equivalence check was retired as a launch blocker after it
        failed, and replaced by the stricter rule that no causal comparison crosses engines.
  \item The activation read-out grid was reduced to the two loci named in the
        pre-registration, dropping four exploratory positions.
  \item The head-level null uses 200 permutations rather than 20.
  \item Maximum sequence length is capped at 4{,}096 tokens for the batched engine, an
        infrastructure limit that cannot affect outputs at these prompt lengths.
  \item Qwen-14B passed the gate, which under the pre-registration triggers a mandatory
        ablation. That run was deferred for compute reasons and is not reported. Its
        predicted outcome under the pattern in Section~\ref{sec:h6res} is a null.
\end{enumerate}
